\subsection{TENSOR PRODUCT OF REPRESENTATIONS}

We have defined in no. 1 the direct sum of a family of representations of g. We shall now define other operations on representations.

Let \(\mathfrak{g}_1, \mathfrak{g}_2\) be two Lie algebras over \(\mathbf{K}\) and \(\mathbf{M}_i\) a \(\mathfrak{g}_i\) -module ( \(i = 1, 2\) ). Let \(\mathbf{U}_i\) be the enveloping algebra of \(\mathfrak{g}_i\) and \(\sigma_i\) the canonical mapping of \(\mathfrak{g}_i\) into \(\mathbf{U}_i\) . Then \(\mathbf{M}_i\) is a left \(\mathbf{U}_i\) -module and hence \(\mathbf{M}_1 \otimes_{\mathbf{K}} \mathbf{M}_2\) has a canonical left \((\mathbf{U}_1 \otimes_{\mathbf{K}} \mathbf{U}_2)\) -module structure. Now \(\mathbf{U}_1 \otimes_{\mathbf{K}} \mathbf{U}_2\) is the enveloping algebra of \(\mathfrak{g}_1 \times \mathfrak{g}_2\) and the mapping \((x_1, x_2) \mapsto \sigma_1(x_1) \otimes 1 + 1 \otimes \sigma_2(x_2)\) is the canonical mapping of \(\mathfrak{g}_1 \times \mathfrak{g}_2\) into this enveloping algebra (§ 2, no. 2). Hence there exists a \((\mathfrak{g}_1 \times \mathfrak{g}_2)\) -module structure on \(\mathbf{M} = \mathbf{M}_1 \otimes_{\mathbf{K}} \mathbf{M}_2\) such that:

\[
\begin{array}{r l} (x _ {1}, x _ {2}) _ {\mathrm{M}}. (m _ {1} \otimes m _ {2}) & = (\sigma_ {1} (x _ {1}) \otimes 1 + 1 \otimes \sigma_ {2} (x _ {2})) . (m _ {1} \otimes m _ {2}) \\ & = ((x _ {1}) _ {\mathrm{M} _ {1}}. m _ {1}) \otimes m _ {2} + m _ {1} \otimes ((x _ {2}) _ {\mathrm{M} _ {2}}. m _ {2}). \end{array}\tag{1}
\]

This structure defines a representation of \(\mathfrak{g}_1\times \mathfrak{g}_2\) on \(\mathbf{M}\) .

If now \(\mathfrak{g}_1 = \mathfrak{g}_2 = \mathfrak{g}\) , the homomorphism \(x \mapsto (x, x)\) of \(\mathfrak{g}\) into \(\mathfrak{g} \times \mathfrak{g}\) , composed with the above representation, defines a representation of \(\mathfrak{g}\) on \(\mathbf{M}\) and hence a \(\mathfrak{g}\) -module structure on \(\mathbf{M}\) such that:

\[
x _ {\mathrm{M}}. (m _ {1} \otimes m _ {2}) = (x _ {\mathrm{M} _ {1}}. m _ {1}) \otimes m _ {2} + m _ {1} \otimes (x _ {\mathrm{M} _ {2}}. m _ {2}).\tag{2}
\]

By an analogous argument we see that:

PROPOSITION 2. Let g be a Lie algebra over K and \(M_{i}\) a g-module ( \(1 \leqslant i \leqslant n\) ). On the tensor product \(M_{1} \otimes_{K} M_{2} \otimes \cdots \otimes M_{n}\) , there exists one and only one g-module structure such that

\[
x _ {\mathrm{M}}. (m _ {1} \otimes \dots \otimes m _ {n}) = \sum_ {i = 1} ^ {n} m _ {1} \otimes \dots \otimes (x _ {\mathrm{M} _ {i}}. m _ {i}) \otimes \dots \otimes m _ {n}\tag{3}
\]

for all \(x \in \mathfrak{g}\) , \(m_1 \in \mathbf{M}_1, \ldots, m_n \in \mathbf{M}_n\) .

The corresponding representation is called the tensor product of the given representations of g on the \(M_{i}\) .

In particular, if \(\mathbf{M}\) is a \(\mathfrak{g}\) -module, Proposition 2 defines a \(\mathfrak{g}\) -module structure
on each \(\mathbf{M}_p = \bigotimes \mathbf{M}\) and hence on the tensor algebra \(\mathbf{T}\) of \(\mathbf{M}\) .

Formula (3) shows that, for all \(x \in \mathfrak{g}\) , \(x_{\mathrm{T}}\) is the unique derivation of the algebra \(\mathrm{T}\) which extends \(x_{\mathrm{M}}\) . We know (§ 2, no. 8) that \(x_{\mathrm{T}}\) defines on passing to the quotient a derivation of the symmetric algebra \(\mathrm{S}\) of \(\mathrm{M}\) . Hence \(\mathrm{S}\) can be considered as a quotient \(\mathfrak{g}\) -module of \(\mathrm{T}\) and the \(x_{\mathrm{S}}\) are derivations of \(\mathrm{S}\) .

Still more particularly, consider g as a g-module by means of the adjoint representation of g. Let U be the enveloping algebra of g. By Proposition 7 of §2, \(x_{\mathrm{M}}\) defines on passing to the quotients a derivation of U which is just the inner derivation defined by \(\sigma(x)\) ( \(\sigma\) denoting the canonical mapping of g into U). Then U can be considered as a quotient g-module of T. If K is a field of characteristic 0, the canonical isomorphism of S onto U is a g-module isomorphism (§2, no. 8).

